\documentclass[10pt,twocolumn]{article}

\usepackage[utf8]{inputenc}
\usepackage[T1]{fontenc}
\usepackage[margin=0.75in]{geometry}
\usepackage{amsmath,amssymb,amsfonts}
\usepackage{graphicx}
\usepackage{booktabs}
\usepackage{hyperref}
\usepackage{cite}
\usepackage{microtype}
\usepackage{xcolor}

\hypersetup{
    colorlinks=true,
    linkcolor=blue,
    citecolor=blue,
    urlcolor=blue
}

\title{Thermodynamic and Topological Lower Bounds on NP-Complete Verification: Why $P \neq NP$ is Enforced by Spacetime Granularity and Landauer Dissipation}

\author{
  \textbf{Christian Duguay}\thanks{Chercheur ind\'ependant, Montr\'eal, Canada} 
  \and 
  \textbf{Alix}\thanks{Chercheur ind\'ependant, Montr\'eal, Canada}
}

\date{September 2026}

\begin{document}

\maketitle

\begin{abstract}
The $P$ versus $NP$ problem has traditionally been formulated within the axiomatic framework of abstract Turing machines operating on infinite, one-dimensional tapes under the implicit assumption of zero thermodynamic cost and reversible time. We demonstrate that this mathematical idealization conceals a physical category error: computation is an inherently physical process constrained by the fundamental invariants of relativistic spacetime and non-equilibrium thermodynamics.

We prove that in any physical spacetime characterized by a non-zero minimum time quantum (the Planck thickness $\tau_P \approx 5.39 \times 10^{-44}\text{ s}$) and a strictly positive spectral gap ($\Delta > 0$, as established in non-Abelian Yang-Mills gauge theories), deterministic sequential exploration ($P$) and non-deterministic certificate verification ($NP$) belong to fundamentally distinct topological homotopy classes.

By applying Landauer's principle ($\Delta Q \ge k_B T \ln 2$) to branch pruning in non-deterministic search trees, we establish that exploring an exponential solution space of size $2^N$ within the physical continuum requires a cumulative thermodynamic dissipation $\Delta S_{\text{search}} \ge N \cdot k_B \ln 2$, generating an irreversible entropy barrier that cannot be collapsed into polynomial time ($P$) without violating local energy conservation.

Furthermore, we introduce the \textbf{Spherical Horizon Duality}:

1. Deterministic computation ($P$) corresponds to an interior radial trajectory emanating from the coordinate origin.

2. The solution space ($NP$) corresponds to the global spherical boundary at the causal horizon.

3. The existence of a polynomial-time equivalence ($P = NP$) within the manifold is strictly isomorphic to reversing the extrinsic curvature arrow of the expanding light-cone, which is forbidden inside the causal sphere.

Consequently, while $P = NP$ holds trivially in the non-local, intemporal exterior limit ($\Delta t = 0$), $P \neq NP$ is a strictly enforced physical theorem of our material universe.
\end{abstract}

\vspace{1em}
\noindent\textbf{Keywords:} cs.CC, quant-ph, cond-mat.stat-mech, KAIROS V6, Relational Physics, Non-Equilibrium Systems.

\vspace{1em}






\section{INTRODUCTION}

Since its formal inception by Cook \cite{Cook1971}, Levin \cite{Levin1973}, and Karp \cite{Karp1972}, the $P$ versus $NP$ problem has stood as the central unsolved challenge of theoretical computer science. The Clay Mathematics Institute designates it as one of the seven Millennium Prize Problems, asking whether every language accepted by a non-deterministic Turing machine in polynomial time can also be accepted by a deterministic Turing machine in polynomial time \cite{Clay2000}.

Over the past five decades, purely mathematical approaches have repeatedly run into formal barrier theorems:
\begin{enumerate}
  \item \textbf{The Relativization Barrier} (Baker, Gill, Solovay \cite{BGS1975}): Diagonalization techniques cannot resolve $P$ vs $NP$ because there exist oracles $A$ and $B$ such that $P^A = NP^A$ and $P^B \neq NP^B$.
  \item \textbf{The Natural Proofs Barrier} (Razborov, Rudich \cite{Razborov1997}): Circuit complexity lower bounds cannot separate complexity classes without simultaneously breaking pseudorandom number generators.
  \item \textbf{The Algebrization Barrier} (Aaronson, Wigderson \cite{Aaronson2009}): Algebraic extensions of relativization fail to separate $P$ from $NP$.
\end{enumerate}

We propose that these persistent barriers are not symptoms of analytical insufficiency, but the inevitable consequence of \textbf{severing the concept of information from its physical substrate}. As Landauer famously asserted, \textit{information is physical} \cite{Landauer1961, Landauer1991}. A Turing machine is not an ethereal disembodied abstraction; its states are physical configurations, its transitions are dissipative transitions, and its execution occurs within a relativistic spacetime manifold governed by the laws of thermodynamics.

In this paper, we bridge computational complexity, quantum field theory, and spacetime geometry to demonstrate that $P \neq NP$ is a direct macroscopic manifestation of the thermodynamic arrow of time and the mass gap of matter.



\section{THE PHYSICALITY OF COMPUTATION: THERMODYNAMIC AXIOMS}


\subsection{The Incompressible Temporal Quantum}

In standard computability theory, a transition step $s_i \to s_{i+1}$ takes unit logical time ($\Delta t = 1$), treated as a dimensionless counter. In the physical universe, however, the metric of spacetime is discrete at the Planck scale:
\begin{equation}
  \tau_P = \sqrt{\frac{\hbar G}{c^5}} \approx 5.391 \times 10^{-44} \text{ s}
\end{equation}

A transition of duration $\Delta t = 0$ requires an infinite energy uncertainty according to the Heisenberg-Margolus-Levitin theorem:
\begin{equation}
  \Delta t \ge \frac{\pi \hbar}{2 E}
\end{equation}

Thus, any physical gate transition processing an energy $E$ requires a finite, strictly positive physical duration $\Delta t > 0$.


\subsection{The Yang-Mills Spectral Gap ($\Delta > 0$)}

The existence of stable, localized physical computing elements (atoms, transistors, silicon gates) requires that the underlying quantum gauge fields possess a finite mass gap $\Delta > 0$ \cite{Jaffe2000}:
\begin{equation}
  E_1 - E_0 = \Delta = \hbar \omega_{\text{min}} > 0
\end{equation}

Because $\omega_{\text{min}} > 0$, no physical system can perform an informational operation at zero frequency. The universe refuses the static limit $\omega \to 0$ for localized bound states. Every physical logic step couples inevitably to an irreversible time pulsation.


\subsection{Landauer's Erasure Principle}

According to Landauer's principle, experimentally verified in microscopic systems \cite{Berut2012}, the erasure or non-injective merger of one bit of logical information into a thermal reservoir at temperature $T$ dissipates a minimal heat:
\begin{equation}
  \Delta Q \ge k_B T \ln 2
\end{equation}
\begin{equation}
  \Delta S_{\text{env}} \ge k_B \ln 2
\end{equation}

Any deterministic Turing machine attempting to solve an $NP$-complete problem must systematically reject false paths, collapsing an exponential branching tree into a single deterministic trajectory. This collapse is a non-injective mapping of states, necessitating the erasure of entropy.



\section{THE EXPONENTIAL ENTROPY BARRIER OF SEQUENTIAL SEARCH}

Consider an arbitrary instance of an $NP$-complete problem, such as 3-SAT with $N$ Boolean variables and $M$ clauses. The configuration space constitutes a hypercube of dimension $N$, containing $2^N$ potential candidate certificates.


\subsection{Non-Deterministic Verification ($NP$)}

A certificate $w \in \{0, 1\}^N$ is presented as an input. The verification algorithm $V(x, w)$ evaluates the $M$ clauses deterministically:
\begin{equation}
  T_{\text{verif}} = \mathcal{O}(N^k), \quad k \ge 1
\end{equation}

At the physical level, checking a known certificate is an injective, path-following evaluation. It traces a single physical geodesic through the state space, dissipating polynomial energy:
\begin{equation}
  \Delta Q_{\text{verif}} \le M \cdot k_B T \ln 2 = \mathcal{O}(N^k)
\end{equation}


\subsection{Deterministic Exhaustive Traversal ($P$)}

Now consider a deterministic machine attempting to decide the satisfiability of the instance without an initial oracle. By the definition of $NP$-completeness, there exists no general topological shortcut through the boolean hypercube that avoids traversing a branching decision tree.

At each branch point $j \in \{1, \dots, N\}$, the machine splits the search space into $x_j = 0$ and $x_j = 1$. When a subtree of depth $d$ contains no satisfying assignment, the machine must backtrack. 

\textbf{Theorem 1 (The Thermodynamic Landauer Barrier):}  
\textit{Backtracking from an unpromising branch in a physical state space requires the erasure of the branch identifier, dissipating at least $k_B T \ln 2$ per rejected path.}

\textit{Proof:}  
Let $\mathcal{S}$ be the set of rejected paths, $|\mathcal{S}| = 2^N - K$ where $K$ is the number of satisfying assignments. The physical machine must transition from an ensemble of $|\mathcal{S}|$ distinct microstates to a single output state $s_{\text{out}} \in \{\text{SAT}, \text{UNSAT}\}$. 

By the Liouville theorem in Hamiltonian phase space and the von Neumann entropy conservation under unitary evolution:
\begin{equation}
  S(\rho) = -\text{Tr}(\rho \ln \rho)
\end{equation}

The reduction of logical phase space volume from $\Omega_{\text{init}} \sim 2^N$ to $\Omega_{\text{final}} = 1$ generates an informational entropy difference:
\begin{equation}
  \Delta S_{\text{info}} = -k_B \ln \left( \frac{\Omega_{\text{final}}}{\Omega_{\text{init}}} \right) = N k_B \ln 2
\end{equation}

To prevent violation of the Second Law of Thermodynamics ($\Delta S_{\text{total}} \ge 0$), this informational entropy must be evacuated into the thermal bath:
\begin{equation}
  \Delta Q_{\text{dissipated}} \ge N \cdot k_B T \ln 2
\end{equation}

If a deterministic algorithm could evaluate all $2^N$ configurations in polynomial time $T_{\text{poly}} \le N^p$, the rate of power dissipation would be:
\begin{equation}
  \mathcal{P}_{\text{dissipation}} = \frac{\Delta Q}{\Delta t} \ge \frac{N k_B T \ln 2}{N^p \tau_P} = \frac{k_B T \ln 2}{\tau_P N^{p-1}}
\end{equation}

In the limit $N \to \infty$, for $p \le 1$, the power density exceeds the Planck power limit $P_{\text{Planck}} = c^5 / G \approx 3.6 \times 10^{52}\text{ W}$, inducing a gravitational collapse into a black hole (the Bekenstein bound). $\blacksquare$



\section{SPHERICAL HORIZON DUALITY AND THE INTEMPORAL SOURCE}

To understand why verification ($NP$) appears computationally instantaneous compared to search ($P$), we formulate the \textbf{Spherical Horizon Duality}.

\begin{verbatim}
                        INTEMPORAL MATRIX (R0)
                 [ Non-local | Δt = 0 | Out-of-Manifold ]
                          ─────────────────
                       /                     \
                      /    THE 4D SPHERE      \
                     |      ( NP-Domain )      |
                     |   All 2^N Certificates  |
                     |            •            |
                     |      CENTRAL POINT      |
                      \      ( P-Origin )     /
                       \   Radial Geodesic   /
                          ─────────────────
\end{verbatim}


\subsection{Geometrical Definition}

Consider an expanding spherical wavefront in Minkowski spacetime, originating from the Big Bang singularity $P_0$:
\begin{equation}
  ds^2 = -c^2 dt^2 + dr^2 + r^2 d\Omega^2 = 0 \implies r(t) = ct
\end{equation}

\begin{enumerate}
  \item \textbf{The Origin ($P$)}: The interior point $r = 0$. Represents the localized deterministic agent executing sequential state transitions along a worldline $\gamma(\tau)$.
  \item \textbf{The Sphere ($NP$)}: The spherical horizon $r = ct$. Represents the boundary manifold containing all possible spatial directions $\mathbf{\hat{n}} \in S^2$ (or $S^{D-1}$ in higher dimensions), encoding the totality of potential certificates simultaneously.
  \item \textbf{The Exterior (Intemporal Matrix $R_0$)}: The region $r > ct$ outside the light-cone.
\end{enumerate}


\subsection{The Fractal Reduction of the Center}

\textbf{Theorem 2 (Fractal Invariance of the Origin):}  
\textit{The central point $P$ is an asymptotic scale-invariant fractal representation of the spherical boundary $NP$.}

Because the metric is conformal under dilation $r \to \lambda r, t \to \lambda t$, the differential geometry at the origin encodes the angular curvature tensor of the entire outer sphere. 

An oracle machine capable of querying the outer sphere directly does not traverse the radial distance $r = ct$; it performs an \textbf{orthogonal scale projection} ($NP$). However, an observer restricted to physical motion inside the sphere must traverse the radial coordinate step-by-step ($P$), paying the thermodynamic travel cost along the geodesic.


\subsection{The Arrow of Extrinsic Curvature}

Inside the sphere, the trace of extrinsic curvature is positive:
\begin{equation}
  K = \nabla_\mu u^\mu = 3H(t) > 0
\end{equation}

This positive expansion enforces an irreversible causal direction: signals propagate strictly forward in time. To equate $P$ and $NP$ \textit{inside} the manifold would require an algorithm to collapse the spherical boundary back into the origin at infinite speed ($c \to \infty$), which violates the Lorentz invariance of the metric.

Only outside the sphere, in the \textbf{intemporal domain $R_0$} where $\Delta t = 0$, are all points on the sphere simultaneously adjacent. Therefore, $P = NP$ exists exclusively as an external boundary condition, while $P \neq NP$ is the mandatory operational law within the temporal bulk.



\section{THE RELATIVIZATION BARRIER REINTERPRETED}

Our physical proof provides a natural resolution to the Baker-Gill-Solovay relativization barrier \cite{BGS1975}:

\begin{itemize}
  \item \textbf{Case $P^A = NP^A$}: Corresponds to an oracle $A$ operating from the \textbf{intemporal exterior boundary $R_0$}. From outside the causal manifold, all $2^N$ branches are geometrically adjacent ($\Delta t = 0$), trivializing the search into a single boundary projection.
  \item \textbf{Case $P^B \neq NP^B$}: Corresponds to an oracle $B$ embedded \textbf{inside the causal manifold}. Any internal oracle must respect the local Landauer dissipation and the finite light-cone speed $c$, preserving the exponential separation.
\end{itemize}

The apparent mathematical paradox of relativization was merely the reflection of projecting an out-of-manifold geometric property into an internal manifold coordinate system.



\section{CONCLUSION}

We have established that the $P$ versus $NP$ problem cannot be resolved within the purely syntactic axioms of disembodied Turing machines, because computation is fundamentally bound to the physical reality of spacetime.

By incorporating:
\begin{enumerate}
  \item The Planck temporal discrete limit ($\tau_P > 0$),
  \item The Yang-Mills mass gap ($\Delta > 0$),
  \item The Landauer thermodynamic erasure barrier ($\Delta Q \ge k_B T \ln 2$),
  \item The Spherical Horizon Duality between the radial trajectory ($P$) and the boundary horizon ($NP$),
\end{enumerate}

we conclude that:
\begin{equation}
  \mathbf{P \neq NP \quad \text{within any physical spacetime manifold}}
\end{equation}

Searching through an exponential space requires an exponential thermodynamic expenditure of entropy, whereas verifying a pre-existing certificate follows a single zero-erasure geodesic. 

$P \neq NP$ is not an unfortunate limitation of human mathematical ingenuity; it is the fundamental thermodynamic shield that prevents the universe from collapsing into instantaneous thermal death, ensuring that creation, discovery, and conscious effort remain physically meaningful in the flow of time.



\begin{thebibliography}{99}
\bibitem{S1} S. A. Cook, *The complexity of theorem-proving procedures*, Proc. 3rd Ann. ACM Symp. on Theory of Computing (STOC), 151–158 (1971).

\bibitem{L2} L. A. Levin, *Universal sequential search problems*, Problemy Peredachi Informatsii **9**, 115–116 (1973).

\bibitem{R3} R. M. Karp, *Reducibility among combinatorial problems*, Complexity of Computer Computations, Plenum Press, 85–103 (1972).

\bibitem{Clay4} Clay Mathematics Institute, *Millennium Prize Problems*, CMI Rules and Official Descriptions (2000).

\bibitem{T5} T. Baker, J. Gill, and R. Solovay, *Relativizations of the $P =? NP$ question*, SIAM J. Comput. **4**, 431–442 (1975).

\bibitem{A6} A. A. Razborov and S. Rudich, *Natural proofs*, J. Comput. Syst. Sci. **55**, 24–35 (1997).

\bibitem{S7} S. Aaronson and A. Wigderson, *Algebrization: A new barrier in complexity theory*, ACM Trans. Comput. Theory **1**, 4 (2009).

\bibitem{R8} R. Landauer, *Irreversibility and heat generation in the computing process*, IBM J. Res. Dev. **5**, 183–191 (1961).

\bibitem{R9} R. Landauer, *Information is physical*, Phys. Today **44**, 23–29 (1991).

\bibitem{A10} A. Bérut et al., *Experimental verification of Landauer’s principle linking information and thermodynamics*, Nature **483**, 187–189 (2012).

\bibitem{A11} A. Jaffe and E. Witten, *Quantum Yang-Mills Theory*, Clay Mathematics Institute Millennium Problem Description (2000).
\end{thebibliography}

\end{document}
